\documentclass[UserManual.tex]{subfiles}
\begin{document}
\begin{sloppypar}

\section{The {\tt CparameterMap} Class}\label{sec:parametermap}

The User will probably write main programs by referencing members of the {\tt CSmoothMaster}, {\tt CTPO} and {\tt CMCMC} classes. It is common that the user will wish to either extract or set parameter that correspond to the various options inherent to those classes. This is done through an object in each of these classes that stores such parameters. These objects are defined by the {\tt CparameterMap} class, and as the name implies, this class provides map-like functionality. 

For example, if the User wishes to write programs that change options during the running, they can use the parameterMap class, which has an object within the {\tt CTPO} class, {\tt tpo->parmap}. For example, the User could set {\tt LAMBDA=3.0} by adding the following line to the main program:
\vspace*{-8pt}{\tt \begin{Verbatim}[commandchars=\\\{\}]
    .
   tpo->parmap.set("TPO_LAMBDA",3.0);
    .
\end{Verbatim}
}\vspace*{-8pt}
This creates a key-value pair, where the key is the string {\tt TPO\_LAMBDA} and the value is 3.0.
If the User wanted to have the program report the current value, they could add a line, e.g.,
\vspace*{-8pt}{\tt \begin{Verbatim}[commandchars=\\\{\}]
    .
   int NMC=tpo->parmap.getI("TPO_NMC",0);
    .
\end{Verbatim}
}\vspace*{-8pt}
Here the first argument is the key. The second argument is the default value, which will be used if the map does not include the given string.
There are similar objects in the {\tt MCMC} class. The User may have created one of these objects, e.g.,
\vspace*{-8pt}{\tt \begin{Verbatim}[commandchars=\\\{\}]
	NBandSmooth::CTPO *tpo=new CTPO();
	NBandSmooth::CSmoothMaster smoothmaster;
	NBandSmooth::CMCMC mcmc(&smoothmaster);
\end{Verbatim}
}\vspace*{-8pt}
All of these three classes have a {\tt CparameterMap} object. For each class the relevant options files are read and the parameter maps are populated during the instantiation. The User can then set or access members of the class through the objects {\tt parmap} member. For example:
\vspace*{-8pt}{\tt \begin{Verbatim}[commandchars=\\\{\}]
	double Lambda=smoothmaster.parmap->getD("SmoothEmulator_LAMBDA",2.5);
	double stepsize=mcmc.parmap.getD("MCMC_METROPOLIS_STEPSIZE",0.06);
\end{Verbatim}
}\vspace*{-8pt}
The {\tt parameterMap} objects were originally populated by reading the options files, {\tt smooth\_data/Options/*options.txt}. Thus, the options can be accessed through the names given in those files.

The {\tt parameterMap} class is basically an extension of C++ maps. The header file, {\tt \$\{SMOOTH\_HOME\}/software/include/msu\_smoothutils/parametermap.h}, is:
\vspace*{-8pt}{\tt \begin{Verbatim}[commandchars=\\\{\}]
class CparameterMap : public map<string,string>\{
public:
  bool getB(string ,bool);
  int getI(string ,int);
  long long int getLongI(string ,long long int);
  string getS(string ,string);
  double getD(string ,double);
  vector< double > getV(string, string);
  vector< string > getVS(string, string);
  vector< vector< double > > getM(string);
  vector< vector< double > > getM(string, double);
  void set(string, double);
  void set(string, int);
  void set(string, unsigned int);
  void set(string, long long int);
  void set(string, bool);
  void set(string, string);
  void set(string, char*);
  void set(string, vector< double >);
  void set(string, vector< string >);
  void set(string, vector< vector< double > >);
  void ReadParsFromFile(const char *filename);
  void ReadParsFromFile(string filename);
  void PrintPars();
  char *message;
\}
\end{Verbatim}
}\vspace*{-8pt}
The first argument of most of the ``{\tt set}'' methods specifies the key in the map, and the second is the value to be set. For the ``{\tt get}'' methods, the first value is the key and the second is the value to be used should the key not exist in the map. The methods {\tt getB,~getI,~getLongI,~getS} and {\tt getD} return booleans, integers, strings or doubles. The software does have the capability to set and find vectors, but that shouldn't be necessary in this instance. All the options listed for {\tt Smooth Emulator} software are set using this functionality.

\end{sloppypar}
\end{document}
